\section{Introduction}
This lab journal will take a shallow dive in to adaptive signal processing, specifically comparison 
of three common types of implementation, the regular \textit{LMS}(REF) and the \textit{Variable step-size LMS}(REF).
Comparison of their performances in an experimental sense, a delve into their specifics, use cases and a critical discussion.
\\
It is a assumed that the reader has some basic understand of digital filtering and a general grasp of their implementation, 
therefore a minimum amount of effort will be spend on that topic, however characterics regarding these
mentioned types of filters, will be explained and discussed.
\\
Filtering in general, is the art of removing the unwanted or highlighting and improving specific 
feature of input signals. Mainly in the analog domain, a filter would be designed for given
frequency or bandwidth, with little possibility of expanding or changing. Such is the 
same problem when designing FIR(REF) or IIR(REF) filters in the digital domains. 
The lack of ability to adapt the characterics of the filtering gave birth to the aptly named \textit{Adaptive filter}.